% problem-000206 9 有机合成选择性控制
\section*{9. 有机合成选择性控制}
有机合成涉及的选择性主要有化学选择性、区域选择性和⽴体选择性等，是学习的重点。

\subsection*{9-1}
画出以下两个氧化反应主要产物A和B的结构式:

（图：）

\subsection*{9-2}
烯基环氧化合物（烯基与环氧直接相连）是重要的合成前体。如下由消旋底物C为原料制备烯基环氧化合物E的路线具有很好的⾮对映选择性（⾮对映选择性过量值>10:1）随后，E可以在⼄酸的作⽤下转化为内酯G1 和 G2。

（图：）

\subsection*{9-2}
-1 画出产物D的⽴体结构式，并判断哪些属于主要产物（在你所画的结构式上圈出即可）

\subsection*{9-2}
-2 由D到E的转化过程属于哪类重排反应类型？

\subsection*{9-2}
-3 画出关键中间体F的⽴体结构式，并判断产物G1中三个⼿性中⼼的绝对构型。

\subsection*{9-3}
烯基环氧化合物H在Lewis酸(LA)催化下经中间体I和J转化为产物；但不同的Lewis酸会形成不同的产物，如LA为 $\mathrm { B F _ { 3 } { \cdot } E t _ { 2 } O }$ 时主要产物为K；⽽TMSOTf催化下则为 $\mathbf { L } _ { \circ }$

（图：）

\subsection*{9-3}
-1 依据以上信息，画出关键中间体I和J及产物K和L的⽴体结构式（分离后没有得到六元环并六元环的产物，K含氧杂氮杂五元环，L含六元环并七元环）。

\subsection*{9-3}
-2 简述为何在 $\mathrm { B F _ { 3 } { \cdot } E t _ { 2 } O }$ 催化下的产物为 $\mathbf { K } _ { \mathtt { C } }$ 。
